\newpage
\section{Summary}

% \subsection{Goals}
This study followed a design-based research (DBR) approach to develop an XR-based training simulator aimed at reducing the learning barrier of FPV quadcopter piloting. The design process was grounded in the integration of human-AI collaboration, with a focus on supporting the development of motor skills.

The study was guided by three research questions: (1) How can FPV quadcopter flight physics be simulated, (2) How can a human-AI collaboration agent be designed to assist users, and (3) What interaction techniques enable high-fidelity XR prototyping.

% \subsection{Results}
The DBR process consisted of eight iterative design cycles. Five cycles focused on refining the XR interaction system through successive prototyping and evaluation. Two cycles were dedicated to the design and training of the HAC agent, and one final cycle integrated both components into a unified training system. Each iteration informed subsequent refinements in both system behavior and interaction design.

% \subsection{Generalization}
The final design demonstrates that FPV piloting skills can be scaffolded through a combination of simplified control schemes and immersive VR-based simulation. In addition, the HAC agent provides real-time adaptive feedback, which supports the formation of experiential learning conditions and reinforces motor skill acquisition through guided interaction.

From a design perspective, the study suggests that coupling human-AI collaboration with XR interaction can enhance skill acquisition in complex control tasks.

% \subsection{Practical and Design Contributions}
Beyond the implemented system, this study contributes (1) an XR prototyping workflow for iterative interaction design, and (2) a framework for developing human-AI collaboration agents. These contributions extend beyond FPV simulation and may inform the design of future XR-based learning systems that require real-time adaptive support.

% \subsection{Future Research}
Future research should extend the current design by incorporating acrobatic flight modes and explicit concept-transfer mechanisms to support far transfer of learning. For XR interaction design, further investigation into multimodal feedback, particularly auditory and haptic channels, is recommended to enhance immersion.

In terms of human-AI collaboration, future work may explore imitation learning and hybrid agent architectures as alternatives or complements to reinforcement learning, with the aim of improving generalization across both near and far transfer learning scenarios.